\documentclass{deimj}

\usepackage[T1]{fontenc}
\usepackage{lmodern}

% 必要に応じてパッケージを導入して下さい。
%\usepackage[dvipdfm]{graphicx}
%\usepackage{latexsym}
%\usepackage{txfonts}
%\usepackage[fleqn]{amsmath}
%\usepackage[psamsfonts]{amssymb}
%\usepackage[deluxe]{otf}

% 印刷位置調整 %
% 必要に応じて値を変更してください．
\hoffset -10mm % <-- 左に 10mm 移動
\voffset -10mm % <-- 上に 10mm 移動

\newcommand{\AmSLaTeX}{%
 $\mathcal A$\lower.4ex\hbox{$\!\mathcal M\!$}$\mathcal S$-\LaTeX}
\newcommand{\PS}{{\scshape Post\-Script}}
\def\BibTeX{{\rmfamily B\kern-.05em{\scshape i\kern-.025em b}\kern-.08em
 T\kern-.1667em\lower.7ex\hbox{E}\kern-.125em X}}

\jtitle{DEIM Forum Class File}
% \jsubtitle{サブタイトル} % サブタイトルを付けたいときはこの行の先頭の % を取る
\authorlist{%
 \authorentry[nishinosono@steng.n-univ.ac.jp]{西之園 萌絵}{Moe Nishinosono}{UnivN}% 
 \authorentry[saikawa@eng.n-univ.ac.jp]{犀川 創平}{Sohei Saikawa}{UnivN}% 
 \authorentry[kunieda@eng.n-univ.ac.jp]{国枝 桃子}{Momoko Kunieda}{UnivN}% 
 \authorentry[magata@magata-lab.co.jp]{真賀田 四季}{Shiki Magata}{Mlab}% 
}
\affiliate[UnivN]{N大学工学部建築学科\hskip1zw
  〒464--8603 愛知県名古屋市千種区不老町}
 {School of Engineering, N University,
  Furo-cho, Chikusa-ku, Nagoya-shi, Aichi 464--8603, Japan}
\affiliate[Mlab]{真賀田研究所\hskip1zw
  〒444--0416 愛知県西尾市一色町妃真加島1番1号}
 {Magata Laboratory,
  1--1 Himakajima, Isshiki-cho, Nishio-shi, Aichi 444--0416, Japan}

% 同一組織で、複数のメールアドレスのドメインがある場合、
% ダガーがずれてしまう。
% その場合、以下のコメントアウトを解除すると、
% 自由にダガーを設定可能であることが分かる。
% \MailAddress{$\dagger$ nishinosono@steng.n-univ.ac.jp,
% \{saikawa,kunieda\}@eng.n-univ.ac.jp,
% $\dagger\dagger$magata@magata-lab.co.jp
% }

\begin{document}
\pagestyle{empty}

\begin{jabstract}
DEIM Forum 論文フォーマット．
\end{jabstract}

\begin{jkeyword}
DEIM，フォーマット，論文執筆の注意事項
\end{jkeyword}

\maketitle

\section{1ページ目のフォーマット}

1ページ目上部には，タイトルと，著者全員についての情報，すなわち，氏名，所属，住所，メールアドレスを記載して下さい．
さらに，あらまし300字程度とキーワードをそれぞれ記述してください．
論文番号は，著者が指定する必要はありません．

\section{原稿提出枚数}

ショートペーパーは2ページ以上，4ページ程度で作成してください．
ロングペーパーは4ページ以上，8ページ程度で原稿を作成してください．

\section{原稿の書き方}

原稿のスタイルは，A4サイズで，9ポイントのフォントを使用し，2段組み，シングルスペースとして下さい．
参考文献は最後に列挙して下さい~\cite{Codd1970}．

% \section*{謝辞}
% 謝辞を記載する場合はこちら．

%\vspace{30mm} <- 文献が本文と近すぎるときは適宜利用してください．
\vspace{2em}

\begin{thebibliography}{99}
\bibitem{Codd1970}
  E. F. Codd, 
  ``A Relational Model of Data for Large Shared Data Banks,''
  Communications of the {ACM} (CACM), Vol. 13, No. 6, pp. 377--387, 1970.
\end{thebibliography}

% bibtexを利用する場合の例
% \bibliographystyle{jplain}
% \bibliography{reference}

\end{document}
